% ============================
% NOTATION & MACROS (IMPROVED SYMBOLS)
% ============================

% ----------------------------
% Global time axis and gaps
% ----------------------------
\newcommand{\TimeAxis}{\mathcal{T}}                        % Global time axis
\newcommand{\timestep}[1]{t_{#1}}                          % Timestamp at index i
\newcommand{\TimeGaps}{\boldsymbol{\Delta}}                % Global time-gap sequence
\newcommand{\timegap}[1]{\delta_{#1}}                      % Time gap at index i

% ----------------------------
% Observations and sequences
% ----------------------------
\newcommand{\obsvec}[1]{\mathbf{x}_{#1}}                   % Observation vector at index i
\newcommand{\ObsSeq}{\mathbf{X}}                           % Sequence matrix
\newcommand{\obsval}[2]{x_{#1#2}}                          % Scalar observation value (i,j)
\newcommand{\obsdim}{D}                                   % Number of observed variables
\newcommand{\seqlen}{L}                                   % Sequence length

% ----------------------------
% Masks and time matrices
% ----------------------------
\newcommand{\Mask}{\mathsf{M}}                             % Mask matrix (sans serif for logical structure)
\newcommand{\mask}[2]{m_{#1#2}}                            % Mask entry (i,j)
\newcommand{\DeltaT}{\mathbf{\Delta t}}                    % Matrix of per-variable time gaps
\newcommand{\deltat}[2]{\Delta t_{#1#2}}                   % Time since last obs of variable j at i
\newcommand{\TimeMat}{\mathbf{T}}                          % Matrix of timestamps

% ----------------------------
% Hidden states
% ----------------------------
\newcommand{\hstate}[1]{\mathbf{h}_{#1}}                   % Hidden state at index i (vector)
\newcommand{\hidden}{\mathbf{h}}                           % Generic hidden state
\newcommand{\hcont}[1]{\mathbf{h}\!\left(#1\right)}         % Continuous-time hidden state h(t)
\newcommand{\cstate}[1]{\mathbf{c}_{#1}}                   % LSTM cell state

% ----------------------------
% Embeddings
% ----------------------------
\newcommand{\embedvec}[1]{\mathbf{e}_{#1}}                 % Embedding vector at index i
\newcommand{\EmbedMat}{\mathbf{E}}                         % Embedding matrix
\newcommand{\Wembed}{\mathbf{W}_{\mathrm{in}}}             % Input projection matrix
\newcommand{\PEmat}{\mathbf{P}}                            % Positional/time embedding matrix

% ----------------------------
% Model weights / projections
% ----------------------------
\newcommand{\Weight}{\mathbf{W}}                           % Generic weight matrix
\newcommand{\Wout}{\mathbf{W}_{\mathrm{out}}}              % Output projection matrix

% ----------------------------
% Attention-specific matrices
% ----------------------------
\newcommand{\Qmat}{\mathbf{Q}}                             % Query matrix
\newcommand{\Kmat}{\mathbf{K}}                             % Key matrix
\newcommand{\Vmat}{\mathbf{V}}                             % Value matrix

% Projection weight matrices for attention
\newcommand{\WQ}{\mathbf{W}_{\!Q}}                          % Projection for queries
\newcommand{\WK}{\mathbf{W}_{\!K}}                          % Projection for keys
\newcommand{\WV}{\mathbf{W}_{\!V}}                          % Projection for values

% Dimensionality of key vectors
\newcommand{\dkey}{d_k}                                   % Key dimension (scalar)

% ----------------------------
% Kernels and convolutions
% ----------------------------
\newcommand{\kernelsize}{\mathcal{K}}                      % Kernel size
\newcommand{\kernelweight}[1]{w_{#1}}                      % Kernel weight

% ----------------------------
% Functions for dynamics
% ----------------------------
\newcommand{\DriftNet}{f_\theta}                          % Drift term network
\newcommand{\DiffNet}{g_\theta}                           % Diffusion term network

% Wiener process
\newcommand{\Wiener}[1]{\mathcal{W}_{#1}}                 % Wiener process at time t

% Data distributions
\newcommand{\pdata}{p_{\mathrm{data}}}                     % Data distribution
\newcommand{\pz}{p_z}                                     % Latent prior distribution

% ----------------------------
% Loss terms
% ----------------------------
\newcommand{\Loss}{\mathcal{L}}                            % Generic loss
\newcommand{\LossRec}{\mathcal{L}_{\mathrm{rec}}}          % Reconstruction loss
\newcommand{\LossPred}{\mathcal{L}_{\mathrm{pred}}}        % Prediction loss
\newcommand{\LossAux}[1]{\mathcal{L}_{\mathrm{aux},#1}}    % Auxiliary loss
\newcommand{\LossTPP}{\mathcal{L}_{\mathrm{TPP}}}          % Temporal Point Process loss
\newcommand{\LossELBO}{\mathcal{L}_{\mathrm{ELBO}}}        % ELBO loss
\newcommand{\LossGAN}{\mathcal{L}_{\mathrm{GAN}}}          % GAN loss

% ----------------------------
% Other math
% ----------------------------
\newcommand{\Intensity}{\lambda}                          % Intensity function
\newcommand{\KL}{\mathrm{KL}}                             % Kullback–Leibler divergence
\newcommand{\Expect}{\mathbb{E}}                          % Expectation

% ============================
% EQUATION MACROS
% ============================

% Global time axis
\newcommand{\EqGlobalTime}{
\begin{equation}
\TimeAxis = ( \timestep{1}, \timestep{2}, \dots, \timestep{\seqlen} ), 
\quad \TimeAxis \in \mathbb{R}^\seqlen .
\label{eq:globaltime}
\end{equation}
}

% Global time gaps
\newcommand{\EqGlobalDelta}{
\begin{equation}
\TimeGaps = ( \timegap{1}, \timegap{2}, \dots, \timegap{\seqlen} ), 
\quad \TimeGaps \in \mathbb{R}^\seqlen .
\label{eq:globaldelta}
\end{equation}
}

% Sequence matrix
\newcommand{\EqSequence}{
\begin{equation}
\ObsSeq = \left( \obsvec{1}, \obsvec{2}, \dots, \obsvec{\seqlen} \right), 
\quad \ObsSeq \in \mathbb{R}^{\obsdim \times \seqlen}.
\label{eq:sequence}
\end{equation}
}

% Mask matrix
\newcommand{\EqMask}{
\begin{equation}
\Mask \in \{0,1\}^{\obsdim \times \seqlen}, \quad
\mask{i}{j} =
\begin{cases}
1 & \text{if variable $x_{ij}$ is observed}, \\
0 & \text{otherwise}.
\end{cases}
\label{eq:mask}
\end{equation}
}

% Δt matrix
\newcommand{\EqDeltaT}{
\begin{equation}
\DeltaT \in \mathbb{R}^{\obsdim \times \seqlen}, \quad  
\deltat{i}{j} = \timestep{i} - t'_{ij},
\label{eq:dt}
\end{equation}
}

% Timestamp matrix
\newcommand{\EqTimeMat}{
\begin{equation}
\TimeMat \in \mathbb{R}^{\obsdim \times \seqlen}, 
\label{eq:T}
\end{equation}
}

% Gaussian Process interpolation
\newcommand{\EqGPInterp}{
\begin{equation}
\big[f(\timestep{1}), \dots, f(\timestep{n})\big]^\top 
\sim \mathcal{N}(m, K), 
\quad m_i = m(\timestep{i}), \quad K_{ij} = k(\timestep{i}, \timestep{j}) .
\label{eq:gp_interp}
\end{equation}
}

% Embedding
\newcommand{\EqEmbed}{
\begin{equation}
\embedvec{i} = \Wembed \obsvec{i}, 
\quad 
\EmbedMat = \left[ \embedvec{1}, \dots, \embedvec{\seqlen} \right].
\label{eq:embed}
\end{equation}
}

% ----------------------------
% Time encoding (sinusoidal) for continuous time
% ----------------------------
\newcommand{\TimeEncoding}[2]{z(\timestep{#1})^{(#2)}}   % Time encoding at time t_i, dim k

\newcommand{\EqTimeEncoding}{
\begin{equation}
\TimeEncoding{i}{k} =
\begin{cases}
\cos\!\left(\tfrac{\timestep{i}}{10000^{(k-1)/M}}\right), & k \ \text{odd}, \\[6pt]
\sin\!\left(\tfrac{\timestep{i}}{10000^{k/M}}\right), & k \ \text{even},
\end{cases}
\label{eq:time_encoding}
\end{equation}
}

% Self-attention
\newcommand{\EqAttention}{
\begin{equation}
\mathrm{Attention}(Q, K, V) = \mathrm{softmax}\!\left(\frac{QK^\top}{\sqrt{d_k}}\right)V .
\label{eq:attention}
\end{equation}
}

% RNN update (discrete)
\newcommand{\EqRNN}{
\begin{equation}
\hstate{i} = \mathrm{RNN}(\hstate{i-1}, \obsvec{i}).
\label{eq:rnn}
\end{equation}
}

% Decay of hidden state (continuous)
\newcommand{\EqDecay}{
\begin{equation}
\hcont{t} = \hstate{i} \cdot f(t - \timestep{i}).
\label{eq:decay}
\end{equation}
}

% Inline imputation
\newcommand{\EqImputation}{
\begin{equation}
\hat{x}_i = (1 - m_i)\, f_\theta(\hstate{i-1}) + m_i\, x_i .
\label{eq:imputation}
\end{equation}
}

\newcommand{\EqConv}{
\begin{equation}
\hstate{i} = \sum_{k=0}^{\kernelsize-1} \kernelweight{k}\,\obsval{i-k}{} .
\label{eq:conv}
\end{equation}
}

% ----------------------------
% Augmented supervised loss
% ----------------------------
\newcommand{\EqAugmentedLoss}{
\begin{equation}
\Loss = \LossPred + \sum_i \lambda_i \LossAux{i} .
\label{eq:augmented_loss}
\end{equation}
}

% Reconstruction loss
\newcommand{\EqRecLoss}{
\begin{equation}
\LossRec = 
\sum_{i \in \mathcal{O}} \| x_i - \hat{x}_i \|^2 .
\label{eq:reconstruction_loss}
\end{equation}
}

% ELBO
\newcommand{\EqELBO}{
\begin{equation}
\LossELBO = 
\Expect_{q_\phi(z|x)} [\log p_\theta(x|z)] 
- \KL\!\left(q_\phi(z|x) \,\|\, p(z)\right).
\label{eq:elbo}
\end{equation}
}

% GAN loss
\newcommand{\EqGAN}{
\begin{equation}
\min_G \max_D \ 
\Expect_{x \sim \pdata} [\log D(x)] 
+ \Expect_{z \sim \pz} [\log (1 - D(G(z)))] .
\label{eq:gan_loss}
\end{equation}
}

% TPP loss
\newcommand{\EqTPP}{
\begin{equation}
\LossTPP = 
\sum_{i=1}^N \log \Intensity(\timestep{i}) 
- \int_0^T \Intensity(s)\,ds .
\label{eq:tpp_loss}
\end{equation}
}

% Hawkes process intensity
\newcommand{\EqHawkes}{
\begin{equation}
\Intensity(t) = \mu + \sum_{\timestep{i} < t} g(t - \timestep{i}) .
\label{eq:hawkes}
\end{equation}
}

% Neural ODE
\newcommand{\EqNODE}{
\begin{equation}
\hcont{t} = \hcont{t_0} + \int_{t_0}^t f_\theta(\hcont{s}, s)\, ds .
\label{eq:node_integral}
\end{equation}
}

% ODE-RNN
\newcommand{\EqODERNN}{
\begin{equation}
\hcont{t} = 
\begin{cases}
\displaystyle \hcont{t_{i-1}} + \int_{t_{i-1}}^{t} f_\theta(\hcont{s}, s)\,ds, & t \in (t_{i-1}, t_i), \\[10pt]
\mathrm{RNN}(\hcont{t}, \obsvec{i}), & t = t_i.
\end{cases}
\label{eq:ode_rnn}
\end{equation}
}

% Neural flow trajectory macro
\newcommand{\NeuralFlow}[2]{F_\theta\bigl(#1, #2\bigr)}

% Neural CDE
\newcommand{\EqNCDE}{
\begin{equation}
\hcont{t} = \hcont{t_0} + \int_{t_0}^t f_\theta(\hcont{s}) \, d\ObsSeq(s) .
\label{eq:ncde}
\end{equation}
}

% SDE
\newcommand{\EqNSDE}{
\begin{equation}
\hcont{t} = \hcont{t_0} 
+ \int_{t_0}^t \DriftNet(\hcont{s}, s)\,ds 
+ \int_{t_0}^t \DiffNet(\hcont{s}, s)\,d\Wiener{s} .
\label{eq:nsde_integral}
\end{equation}
}

